\documentclass[aspectratio=169,10pt]{beamer}
\usetheme[numbering=fraction,progressbar=frametitle]{metropolis}
\usepackage{fontspec}
\setsansfont{Latin Modern Sans}
\usepackage{amsmath,amssymb,mathtools,graphicx,booktabs}
\usepackage{polyglossia}
\setdefaultlanguage{french}
\setbeamertemplate{navigation symbols}{}
\setbeamerfont{frametitle}{size=\large}
\setbeamerfont{normal text}{size=\normalsize}
\setlength{\parskip}{.45em}
\hypersetup{pdfauthor={},pdfsubject={Lecture slides},pdfkeywords={},pdftitle={Problèmes de synthèse, raisonnement et automatismes}}
\providecommand{\N}{\mathbb{N}}
\providecommand{\Z}{\mathbb{Z}}
\providecommand{\D}{\mathbb{D}}
\providecommand{\Q}{\mathbb{Q}}
\providecommand{\R}{\mathbb{R}}
\providecommand{\C}{\mathbb{C}}
\providecommand{\vect}[1]{\overrightarrow{#1}}
\providecommand{\vi}{\vec{\imath}}
\providecommand{\vj}{\vec{\jmath}}
\providecommand{\vk}{\vec{k}}
\providecommand{\norm}[1]{\left\lVert #1\right\rVert}
\providecommand{\vc}[2]{\begin{pmatrix}#1\\#2\end{pmatrix}}
\providecommand{\deter}[2]{\det\left(#1,#2\right)}
\providecommand{\abs}[1]{\left\lvert #1\right\rvert}
\providecommand{\dd}{\,\mathrm{d}}
\providecommand{\e}{\mathrm{e}}
\providecommand{\Card}{\operatorname{Card}}
\providecommand{\E}{\operatorname{E}}
\providecommand{\V}{\operatorname{V}}
\providecommand{\degre}{^{\circ}}
\providecommand{\un}[1]{\,\mathrm{#1}}
\providecommand{\sci}[2]{#1\times10^{#2}}
\providecommand{\pgcd}{\operatorname{PGCD}}
\providecommand{\ppcm}{\operatorname{PPCM}}
\title{Problèmes de synthèse, raisonnement et automatismes}
\subtitle{Cours et exemples}
\date{}
\author{}
\begin{document}
\begin{frame}[plain]
\vfill
{\usebeamerfont{title}\inserttitle\par}
\medskip
{\usebeamerfont{subtitle}\insertsubtitle\par}
\medskip
{\usebeamercolor[fg]{progress bar}\rule{\textwidth}{.5pt}}
\vfill
\end{frame}
\begin{frame}{Implication, réciproque et contre-exemple}
« Si \(P\), alors \(Q\) » est une implication \(P\Rightarrow Q\). Sa réciproque est \(Q\Rightarrow P\) ; elle n'est pas automatiquement vraie. Une équivalence \(P\iff Q\) exige les deux implications.
Un contre-exemple suffit à réfuter une assertion universelle ; quelques exemples ne suffisent pas à la prouver. La contraposée \(\text{non }Q\Rightarrow\text{non }P\) est équivalente à \(P\Rightarrow Q\).
\end{frame}
\begin{frame}{Implication, réciproque et contre-exemple}
Pour un entier \(n\) : « Si \(n\) est multiple de \(4\), alors \(n\) est pair. » Démontrer l'implication, écrire sa réciproque et sa contraposée, puis dire lesquelles sont vraies.
\end{frame}
\begin{frame}{Correction : Implication, réciproque et contre-exemple}
Si \(n=4k\), alors \(n=2(2k)\) est pair. Réciproque : « Si \(n\) est pair, alors \(n\) est multiple de \(4\) » ; elle est fausse avec \(n=2\). Contraposée : « Si \(n\) est impair, alors \(n\) n'est pas multiple de \(4\) » ; elle est vraie et équivalente à l'implication initiale.
\end{frame}
\begin{frame}{Implication, réciproque et contre-exemple}
Pour \(x\in\R\), étudier « \(x=3\Rightarrow x^2=9\) » et sa réciproque. Donner une équivalence correcte.
\end{frame}
\begin{frame}{Correction : Implication, réciproque et contre-exemple}
L'implication est vraie. La réciproque est fausse avec \(x=-3\). L'équivalence correcte est \(x^2=9\iff x=3\text{ ou }x=-3\). Avec l'hypothèse supplémentaire \(x\geq0\), on peut écrire \(x^2=9\iff x=3\).
\end{frame}
\begin{frame}{Négation, disjonction de cas et preuve}
Nier correctement : « \(x>2\) » ; « \(x\in[1;4]\) » ; « Tous les élèves font du sport » ; « \(A\) et \(B\) sont réalisés ».
\end{frame}
\begin{frame}{Correction : Négation, disjonction de cas et preuve}
« \(x\leq2\) » ; « \(x<1\) ou \(x>4\) » ; « Au moins un élève ne fait pas de sport » ; « \(A\) n'est pas réalisé ou \(B\) n'est pas réalisé ». Le « ou » est inclusif.
\end{frame}
\begin{frame}{Négation, disjonction de cas et preuve}
Montrer par disjonction de cas que \(n(n+1)\) est pair pour tout entier \(n\). Montrer ensuite que \(n^2+n+1\) est impair.
\end{frame}
\begin{frame}{Correction : Négation, disjonction de cas et preuve}
Si \(n\) est pair, \(n(n+1)\) a un facteur pair. Si \(n\) est impair, \(n+1\) est pair ; le produit est encore pair. Les deux cas couvrent tous les entiers. Comme \(n^2+n+1=n(n+1)+1\), c'est un entier impair.
\end{frame}
\begin{frame}{Négation, disjonction de cas et preuve}
Un élève écrit « \(n^2\) pair implique \(n\) pair, parce que je l'ai vérifié pour \(n=2,4,6\) ». Remplacer cette justification par une preuve.
\end{frame}
\begin{frame}{Correction : Négation, disjonction de cas et preuve}
Si \(n\) était impair, le résultat démontré en séquence 1 donnerait \(n^2\) impair, contradiction. C'est un raisonnement par contraposée, ou présenté ici par l'absurde. Trois exemples ne traitent pas tous les entiers.
\end{frame}
\begin{frame}{Croiser algèbre, fonctions et pourcentages}
Deux tarifs fictifs de location sont proposés pour \(x\) heures, avec \(x\in[0;8]\) : \(A(x)=3000+2000x\) Ar et \(B(x)=3500x\) Ar.

\begin{enumerate}
\setcounter{enumi}{0}\item Calculer les coûts pour \(1\) h et \(4\) h.
\item Trouver l'égalité des tarifs puis les durées où \(A\) est moins cher.
\item Une remise de \(10\,\%\) est accordée sur tout le tarif \(A\). Donner le nouveau tarif \(A_r\) et le nouveau seuil de comparaison.

\end{enumerate}
\end{frame}
\begin{frame}{Correction : Croiser algèbre, fonctions et pourcentages}
1. À \(1\) h : \(5000\) Ar contre \(3500\) Ar. À \(4\) h : \(11000\) Ar contre \(14000\) Ar. 2. \(3000+2000x=3500x\iff x=2\). \(A<B\) pour \(x>2\), donc sur \(]2;8]\). 3. \(A_r(x)=0{,}9(3000+2000x)=2700+1800x\). L'égalité avec \(B\) donne \(1700x=2700\), soit \(x=27/17\approx1{,}59\) h. \(A_r<B\) sur \(]27/17;8]\). La remise s'applique aussi au forfait fixe.
\end{frame}
\begin{frame}[fragile]{Python}
\small
\begin{verbatim}
def tarif_moins_cher(x):
    a = 3000 + 2000*x
    b = 3500*x
    if a < b:
        return "A"
    elif b < a:
        return "B"
    else:
        return "egalite"
\end{verbatim}
\end{frame}
\begin{frame}{Croiser algèbre, fonctions et pourcentages}
Prévoir les résultats pour \(x=1,2,4\). Modifier la ligne qui calcule \(a\) pour prendre en compte la remise.
\end{frame}
\begin{frame}{Correction : Croiser algèbre, fonctions et pourcentages}
Résultats : \texttt{"B"}, \texttt{"egalite"}, \texttt{"A"}. Remplacer par \texttt{a = 0.9*(3000 + 2000*x)}. Le domaine \([0;8]\) est supposé respecté lors des appels.
\end{frame}
\end{document}
